\documentclass[profile=standard]{ipara-handbook}
\title{B00｜内积、正交与投影}
\author{}
\date{}
\begin{document}
\maketitle

\begin{quote}
这座桥回答一个最基础的跨科问题：\textbf{空间几何里的“垂直、分量、最近”，线代里的“正交、坐标”，以及 Fourier 里的“系数”，为什么能够使用同一种语言？}
\end{quote}

\section{本册定位：在线性结构上增加一套几何读法}
线代 Topic01 Own 向量空间、基、坐标、内积、正交与 Gram--Schmidt；高数 Topic06 Own 空间对象、方向、角度、距离等几何计算；高数 Topic11 Own Fourier 展开与收敛。B00 不重新教授这些本体，而只 Own 一条稳定翻译：
\[
\boxed{
\text{内积}
\xrightarrow{\text{定义几何}}
\text{长度/角度/正交}
\xrightarrow{\text{正交分解}}
\text{投影 + 正交残差}
\xrightarrow{\text{坐标读取}}
\text{最小误差/正交系数}
}
\]

这里第一个箭头表示“在线性空间上增加额外结构”，并不是向量空间公理自动推出内积。

\section{为什么普通坐标还不够}
普通基只保证每个向量有唯一坐标。若题目进一步问“长度多大、夹角多少、哪个方向最接近、怎样把不同方向解耦”，还需要内积。

在标准欧氏空间中
\[
\langle x,y\rangle=x^Ty,
\qquad
\lVert x\rVert=\sqrt{\langle x,x\rangle}.
\]
正交由
\[
\langle x,y\rangle=0
\]
定义。于是“垂直”不再依赖一张图，而成为可计算的代数关系。

\subsection{坐标点积不是永远的抽象内积公式}
若换到一般非规范正交基，坐标列的普通点积不再自动等于原空间内积。真正的几何对象是内积本身；矩阵坐标必须连同基或 Gram 矩阵一起解释。

因此 B00 的第一个对象边界是
\[
\boxed{\text{内积几何}\neq\text{某一份坐标中的机械点积}.}
\]

\section{母接口：投影 = 保留目标分量，残差与目标正交}
先看一维子空间 $W=\operatorname{span}\{u\}$，$u\neq0$。寻找 $v$ 在 $W$ 中的“最接近分量”，写成 $cu$。要求残差
\[
r=v-cu
\]
与 $u$ 正交：
\[
\langle v-cu,u\rangle=0.
\]
于是
\[
\boxed{
\operatorname{proj}_{u}v
=\frac{\langle v,u\rangle}{\langle u,u\rangle}u
}
\]
并且
\[
v=\operatorname{proj}_{u}v+r,
\qquad r\perp u.
\]

这条正交残差条件不是附加检查，而是投影的接口不变量。

\section{为什么正交投影同时解决“最近点”}
对任意 $w\in W$，设 $p=\operatorname{proj}_Wv$。由于
\[
v-p\perp W,
\]
而 $p-w\in W$，所以两部分正交。由勾股关系
\[
\lVert v-w\rVert^2
=\lVert v-p\rVert^2+\lVert p-w\rVert^2
\ge \lVert v-p\rVert^2.
\]
等号只在 $w=p$ 时成立。因此
\[
\boxed{p\text{ 是 }W\text{ 中距离 }v\text{ 最近的点}.}
\]

这一步解释了为什么“画垂线”“最短距离”“最小平方误差”会在不同章节重复出现：它们共享的是同一条正交分解机制。

\section{多方向：正交基为什么让坐标直接变成投影系数}
若 $e_1,\dots,e_k$ 两两正交，则对 $v\in\operatorname{span}\{e_i\}$，
\[
v=\sum_{i=1}^k c_i e_i,
\qquad
c_i=\frac{\langle v,e_i\rangle}{\langle e_i,e_i\rangle}.
\]
若进一步规范化为 $\lVert e_i\rVert=1$，则
\[
\boxed{c_i=\langle v,e_i\rangle.}
\]

所以正交基的核心收益不是“看起来整齐”，而是每个坐标都可以单独读取，不会被其他方向的分量污染。

这正是 B08 能把 Fourier 系数解释成投影坐标的根源。

\section{母例一：空间几何中的垂足}
取
\[
v=(2,1)^T,
\qquad u=(1,1)^T.
\]
则
\[
p=\operatorname{proj}_u v
=\frac{3}{2}(1,1)^T,
\]
残差
\[
r=v-p=\left(\frac12,-\frac12\right)^T,
\qquad r^Tu=0.
\]
所以 $p$ 不仅是“某个公式算出的向量”，还是直线 $\operatorname{span}\{u\}$ 上离 $v$ 最近的点。

\section{母例二：函数也可以有“正交分量”}
在函数空间中定义
\[
\langle f,g\rangle=\int_{-1}^{1}f(x)g(x)\,dx.
\]
把 $f(x)=x^2$ 投影到常数函数子空间 $W=\operatorname{span}\{1\}$。系数为
\[
c=\frac{\langle x^2,1\rangle}{\langle1,1\rangle}
=\frac{\int_{-1}^{1}x^2\,dx}{\int_{-1}^{1}1\,dx}
=\frac13.
\]
于是
\[
x^2=\frac13+\left(x^2-\frac13\right),
\]
且
\[
\int_{-1}^{1}\left(x^2-\frac13\right)dx=0.
\]
同一个“投影 + 正交残差”接口已经从二维向量无缝迁移到了函数对象。

\section{B00 究竟串起哪些章节}
\begin{tabular}{p{0.23\linewidth}p{0.30\linewidth}p{0.39\linewidth}}
\hline
Owner / 下游 & 提供或接收的对象 & B00 的翻译责任\\
\hline
高数 Topic06 空间对象 & 方向、法向、角度、距离 & 把几何垂直/分量翻译成内积与正交投影\\
线代 Topic01 & 内积、正交基、Gram--Schmidt & 提供抽象代数本体，B00 解释其跨科几何意义\\
B04 约束极值 & 切空间与法空间 & 用正交补解释“目标梯度的切向投影必须为零”\\
高数 Topic11 / B08 & Fourier 基函数与系数 & 把函数系数翻译成正交投影坐标\\
\hline
\end{tabular}

因此 B00 是 B04、B08 的共同上游，而不是一篇孤立的“投影公式笔记”。

\section{资格与停止边界}
\begin{enumerate}
\item 必须先说明使用哪个内积；换内积会改变长度、正交与投影。
\item 投影到一个子空间前，要确认目标集合真的是线性/仿射对象；一般曲线上的最近点不能直接套线性投影公式。
\item 普通基不等于正交基；只有正交基才能逐项用投影系数解耦坐标。
\item B00 只解释投影接口。Fourier 的无限级数收敛与恢复值转 B08/高数 Topic11；约束极值的正则性与候选全集转 B04/高数 Topic07。
\end{enumerate}

\section{反桥梁：最容易越级的四组关系}
\begin{tabular}{p{0.20\linewidth}p{0.04\linewidth}p{0.20\linewidth}p{0.48\linewidth}}
\hline
A & & B & 真正区别\\
\hline
线性无关 & $\neq$ & 正交 & 正交非零组必无关；无关只排除线性冗余，不要求内积为零\\
正交 & $\neq$ & 单位化 & 前者比较方向，后者只规范长度\\
正交 & $\neq$ & 概率独立 & 一个由内积定义，一个由联合分布分解定义\\
投影 & $\neq$ & 随意删坐标 & 正交投影由目标子空间和内积共同决定，残差必须满足正交不变量\\
\hline
\end{tabular}

\section{什么时候应该调用这座桥}
题面出现“最近、最短距离、垂足、正交系数、最小误差、正交基、Fourier 系数来源”时，第一观察不是背公式，而是问：
\begin{enumerate}
\item 当前对象是什么空间里的对象？
\item 使用哪个内积？
\item 目标子空间是什么？
\item 能否写成“目标分量 + 正交残差”？
\end{enumerate}

\section{一页压缩}
B00 最后只保留一句话：
\[
\boxed{
\text{正交投影把对象分成“目标子空间中的可保留分量”与“对目标完全不可见的正交残差”。}
}
\]
忘掉公式时，从
\[
v=p+r,\qquad p\in W,\quad r\in W^\perp
\]
重新推出投影系数；再用勾股关系恢复“最近点”含义。若这个结构能同时解释空间距离、正交坐标与 Fourier 系数，B00 才算真正掌握。

\end{document}
